\documentclass{article}
\usepackage{amsmath}
\usepackage{amssymb}
\usepackage{geometry}
\usepackage{graphicx}
\usepackage{xepersian}
\settextfont{XB Niloofar}

\geometry{a4paper, margin=1in}

\title{\lr{Understanding the Covariance Matrix} \\ (درک ماتریس کوواریانس)}
\author{}
\date{}
\begin{figure}[h!]
  \centering
  \includegraphics[width=0.75\textwidth]{img_page_4_1.png}
  \caption{نقشه حرارتی ماتریس کوواریانس نمونه برای داده‌های شبیه‌سازی‌شده}
\end{figure}

\begin{figure}[h!]
  \centering
  \includegraphics[width=0.75\textwidth]{img_page_5_1.png}
  \caption{نمودار توزیع داده‌ها بر اساس همبستگی و پراکندگی آن‌ها}
\end{figure}

\begin{figure}[h!]
  \centering
  \includegraphics[width=0.75\textwidth]{img_page_7_1.png}
  \caption{رابطه مستقیم و همبستگی مثبت بین دو متغیر تصادفی}
\end{figure}

\begin{figure}[h!]
  \centering
  \includegraphics[width=0.75\textwidth]{img_page_10_1.png}
  \caption{نمایش بردارهای ویژه و مقادیر ویژه به عنوان جهت‌های اصلی واریانس و پراکندگی داده‌ها}
\end{figure}

\begin{figure}[h!]
  \centering
  \includegraphics[width=0.75\textwidth]{img_page_16_1.png}
  \caption{ماتریس همبستگی پیرسون مربوط به شباهت امتیازدهی کاربران به فیلم‌ها}
\end{figure}

\begin{figure}[h!]
  \centering
  \includegraphics[width=0.75\textwidth]{img_page_17_1.png}
  \caption{خلاصه هندسی و مفهومی ماتریس کوواریانس}
\end{figure}














\begin{document}

\maketitle

\section*{فهرست مطالب}
\begin{enumerate}
    \item کوواریانس چیست؟ (\lr{What is Covariance?})
    \item ماتریس کوواریانس چیست؟ (\lr{What is the Covariance Matrix?})
    \item مقادیر ویژه و بردارهای ویژه ماتریس کوواریانس
    \item ارتباط با تحلیل مولفه‌های اصلی (\lr{PCA})
    \item محاسبه دستی ماتریس کوواریانس نمونه
    \item ماتریس همبستگی (\lr{Correlation Matrix})
    \item مثال: فیلتر کردن داده‌ها بر اساس همبستگی پیرسون (\lr{Pearson Correlation})
    \item نتیجه‌گیری
\end{enumerate}

\newpage

\section*{۱. کوواریانس چیست؟}
کوواریانس (\lr{Covariance}) میزان تغییر همزمان دو متغیر تصادفی را اندازه‌گیری می‌کند. اگر متغیرها رفتار مشابهی نشان دهند (یعنی با افزایش یکی، دیگری نیز افزایش یابد)، کوواریانس مثبت است. اگر رفتار متضاد نشان دهند (با افزایش یکی، دیگری کاهش یابد)، کوواریانس منفی است. اگر متغیرها مستقل باشند، کوواریانس صفر است (اگرچه کوواریانس صفر لزوماً به معنای استقلال نیست).

فرمول کوواریانس نمونه بین دو متغیر X و Y به صورت زیر است:
\[
cov(X,Y) = \frac{1}{n-1} \sum_{i=1}^n (x_i - \bar{x})(y_i - \bar{y})
\]

\section*{۲. ماتریس کوواریانس چیست؟}
برای مجموعه‌داده‌ای با چندین متغیر (ویژگی)، ماتریس کوواریانس، کوواریانس‌های بین تمام جفت متغیرها را سازماندهی می‌کند. برای مجموعه‌داده‌ای با $d$ متغیر ($X_1, X_2, \dots, X_d$)، ماتریس کوواریانس $\Sigma$ (یا $S$ برای نمونه) یک ماتریس متقارن $d \times d$ است که در آن:
\begin{itemize}
    \item عنصر در موقعیت $(i, j)$ کوواریانس بین متغیر $X_i$ و متغیر $X_j$ است: $\Sigma_{ij} = cov(X_i, X_j)$.
    \item عناصر روی قطر اصلی $\Sigma_{ii}$ واریانس هر متغیر هستند: $\Sigma_{ii} = cov(X_i, X_i) = var(X_i)$.
\end{itemize}

مثال برای ۳ متغیر $(X, Y, Z)$:
\[
S = \begin{pmatrix} 
var(X) & cov(X,Y) & cov(X,Z) \\ 
cov(Y,X) & var(Y) & cov(Y,Z) \\ 
cov(Z,X) & cov(Z,Y) & var(Z) 
\end{pmatrix}
\]
از آنجا که $cov(X, Y) = cov(Y, X)$ است، این ماتریس همواره متقارن است.

\subsection*{مثال شهودی}
فرض کنید مجموعه‌داده‌ای از ۱۵۰ دانش‌آموز با متغیرهای ساعت مطالعه، ساعت تفریح و نمره ریاضی داریم:
\begin{itemize}
    \item \textbf{واریانس‌ها (عناصر روی قطر):} واریانس ساعت مطالعه (۱۴.۲۴)، واریانس ساعت تفریح (۱۲.۴۹) و واریانس نمره ریاضی (۱۲.۳۶) که پراکندگی هر متغیر را نشان می‌دهند.
    \item \textbf{کوواریانس‌ها (عناصر خارج از قطر):}
    \begin{itemize}
        \item کوواریانس ساعت مطالعه و ساعت تفریح (۱۰.۱۱-): نشان‌دهنده رابطه معکوس است (دانش‌آموزانی که بیشتر مطالعه می‌کنند، کمتر تفریح می‌کنند).
        \item کوواریانس ساعت مطالعه و نمره ریاضی (۸.۲۶+): نشان‌دهنده رابطه مستقیم است.
        \item کوواریانس ساعت تفریح و نمره ریاضی (۶.۶۸-): نشان‌دهنده رابطه معکوس است.
    \end{itemize}
\end{itemize}

\section*{۳. مقادیر ویژه و بردارهای ویژه ماتریس کوواریانس}
بردارهای ویژه و مقادیر ویژه اطلاعات مهمی درباره ساختار داده‌ها ارائه می‌دهند:
\begin{itemize}
    \item \textbf{بردارهای ویژه (\lr{Eigenvectors}):} جهت‌های حداکثر واریانس را در داده‌ها نشان می‌دهند که همان محورهای اصلی توزیع داده‌ها هستند.
    \item \textbf{مقادیر ویژه (\lr{Eigenvalues}):} مقدار واریانس داده‌ها را در امتداد جهت بردارهای ویژه نشان می‌دهند. مقدار ویژه بزرگتر یعنی واریانس بیشتر در آن جهت.
\end{itemize}

\section*{۴. ارتباط با تحلیل مولفه‌های اصلی (\lr{PCA})}
این تحلیل هسته اصلی الگوریتم \lr{PCA} است:
\begin{enumerate}
    \item بردارهای ویژه ماتریس کوواریانس، جهت‌های مولفه‌های اصلی هستند.
    \item مقادیر ویژه نشان‌دهنده میزان واریانس ثبت‌یافته توسط هر مولفه اصلی متناظر هستند.
\end{enumerate}
با تصویر داده‌ها روی بردارهای ویژه، داده‌ها به یک سیستم مختصات جدید منتقل می‌شوند که محورهای آن جهت‌های بیشترین واریانس هستند. این روش برای کاهش ابعاد با حذف مولفه‌های دارای مقادیر ویژه کوچک استفاده می‌شود.

\section*{۵. محاسبه دستی ماتریس کوواریانس نمونه}
فرمول کوواریانس نمونه بدون انحراف بین دو متغیر $X_i$ و $X_j$ به صورت زیر است:
\[
S_{ij} = \frac{1}{n-1} \sum_{k=1}^n (x_{ki} - \bar{x}_i)(x_{kj} - \bar{x}_j)
\]
که در آن $n$ تعداد نمونه‌ها، $x_{ki}$ مشاهده $k$-ام از ویژگی $i$-ام و $\bar{x}_i$ میانگین نمونه ویژگی $i$-ام است.

برای محاسبه ماتریسی:
\begin{enumerate}
    \item \textbf{مرکزی‌سازی داده‌ها:} میانگین هر ستون را از مقادیر آن ستون کم می‌کنیم تا ماتریس داده‌های مرکزی‌شده $X_{\text{centered}}$ به دست آید:
    \[
    X_{\text{centered}} = X - \mu^T
    \]
    \item \textbf{محاسبه ماتریس کوواریانس:} با استفاده از ضرب ماتریسی داریم:
    \[
    S = \frac{1}{n-1} X_{\text{centered}}^T X_{\text{centered}}
    \]
\end{enumerate}
این رویکرد ماتریسی بسیار کارآمد بوده و برای مجموعه‌داده‌های بزرگ ترجیح داده می‌شود.

\section*{۶. ماتریس همبستگی}
هرچند ماتریس کوواریانس جهت و اندازه رابطه خطی را نشان می‌دهد، اما اندازه آن به مقیاس متغیرها بستگی دارد. ماتریس همبستگی (\lr{correlation matrix}) که با $R$ نشان داده می‌شود، کوواریانس را استاندارد می‌کند و قدرت و جهت رابطه خطی را بین جفت متغیرها بین -۱ و ۱ اندازه‌گیری می‌کند:
\[
R_{ij} = \frac{Cov(X_i, X_j)}{\sqrt{Var(X_i) Var(X_j)}} = \frac{S_{ij}}{\sqrt{S_{ii} S_{jj}}}
\]
عناصر روی قطر اصلی همواره ۱ هستند.

\section*{۷. مثال: فیلتر کردن داده‌ها بر اساس همبستگی پیرسون}
ضریب همبستگی پیرسون (\lr{Pearson correlation}) که از ماتریس کوواریانس به دست می‌آید، برای شناسایی و حذف ویژگی‌های تکراری یا ضعیف استفاده می‌شود. به عنوان مثال، در داده‌های امتیازدهی فیلم کاربران، کاربران با بیشترین شباهت ('Robert', 'Clara') دارای همبستگی ۱.۰۰ و کاربران با کمترین شباهت ('Dan', 'Robert') دارای همبستگی -۰.۹۷ هستند.

\section*{۸. نتیجه‌گیری}
ماتریس کوواریانس ابزاری بنیادی در آمار چندمتغیره است.
\begin{itemize}
    \item این ماتریس مربع و متقارن است.
    \item عناصر روی قطر اصلی واریانس و عناصر خارج از قطر کوواریانس هستند.
    \item بردارهای ویژه و مقادیر ویژه آن محورهای اصلی و اندازه پراکندگی داده‌ها را مشخص می‌کنند که مبنای \lr{PCA} است.
\end{itemize}

\end{document}
